\subsection{The birth-weight paradox}
\totalscore{6}

In this exercise we investigate the \emph{birth-weight paradox}.
In his \emph{Book of Why}, Judea Pearl describes this paradox as follows:
\begin{quote}
In the mid-1960's, Jacob Yerushalmy pointed out that a mother's smoking during
pregnancy seemed to benefit the health of her newborn baby, if the baby happened to be born underweight.
This puzzle, called the birth-weight paradox, flew in the face of the emerging medical consensus 
about smoking and was not satisfactorily explained until 2006---more than forty years after
Yerushalmy's original paper.
[...]

In 1959, Yerushalmy had launched a long-term public health study that collected pre- and postnatal
data on more than 15,000 children in the San Francisco Bay Area. The data included information on 
mother's smoking habits, as well as the birth weights and mortality rates of their babies in the first month of life.
Several studies had already shown that the babies of smoking mothers weighed less at birth on average than the babies of
nonsmokers, and it was natural to suppose that this would translate to poorer survival. Indeed, a nationwide study of
low-birth-weight infants had shown that their death rate was more than twenty times higher than that of normal-birth-weight infants.
Thus, epidemiologists posited a chain of causes and effects: smoking causes low birth weight, which causes mortality.
[...]

What Yerushalmy found in the data was unexpected even to him. It was true that the babies of smokers were lighter
on average than the babies of nonsmokers. However, the low-birth-weight babies of smoking mothers had a \emph{better}
survival rate than those of nonsmoking mothers. It was as if the mother's smoking actually had a protective effect.
[...]

Most modern epidemiologists believe that smoking does increase neonatal mortality---for example, because it 
interferes with oxygen transfer across the placenta. But how can we reconcile this hypothesis with the data?
[...]

In fact, Yerushalmy's data were completely consistent with the model that smoking causes low birth weight
which causes mortality, once we add a little more to it. Smoking may be harmful in that it contributes to
low birth weight, but certain other causes of low birth weight, such as serious or life-threatening 
genetic abnormalities, are much more harmful.
% There are two possible explanations for low birth weight in one particular baby: it might have a smoking mother, or it might be affected by one of those other causes.
\end{quote}

\begin{enumerate}[label=\alph*)]
  \item Draw a causal graph that matches with Pearl's story.
    Make sure to include the four observed variables in the table:
    \begin{center}\small\begin{tabular}{lll}
      Symbol & Meaning & Values \\
      \hline
      $S$ & Smoking mother               & 0: does not smoke, 1: smokes \\ 
      $W$ & Birth weight of child        & 0: normal, 1: low \\
      $A$ & Genetic abnormality of child & 0: normal, 1: abnormality \\
      $M$ & Mortality of child           & 0: lives, 1: died \\
    \end{tabular}\end{center}
    Include the minimal number of edges that you need to make it compatible with the story.
    \score{3}
    \answer{
    \begin{tikzpicture}
      \node[ndout] (S) at (0,0) {$S$};
      \node[ndout] (W) at (2,-1) {$W$};
      \node[ndout] (A) at (4,0) {$A$};
      \node[ndout] (M) at (2,-3) {$M$};
      \draw[tuh] (S) to (W);
      \draw[tuh] (S) to (M);
      \draw[tuh] (W) to (M);
      \draw[tuh] (A) to (W);
      \draw[tuh] (A) to (M);
    \end{tikzpicture}
    }

  \item With the help of the causal graph, explain in words how the paradox can be resolved.
    \score{3}
    \answer{
    Let us assume smoking increases the risk of mortality, but not as much as a genetic abnormality does.

    Then this is a classical ``explaining away'' phenomenon.
    While $S$ and $A$ are independent, they become negatively correlated when conditioning on $W$ (low birth-weight).
    There are two possible explanations for low birth weight in one particular baby: it might have a smoking mother, or it might be affected by a genetic abnormality.
    Therefore, for low-birth weight babies: if the mother is found to be non-smoking, it is quite likely a genetic abnormality, strongly increasing the risk of mortality;
    while if the mother is found to be smoking, this already explains the low birth-weight, and the chance of a genetic abnormality stays small (as these are rare in the general population), and the risk of mortality increases, but not as much.
  }
   \points{This can be formalized as follows.
    %Assume $P(W=1 \given \doit(S,A)) = w_0 e^{\omega_S S + \omega_A A}$, which models $S$ and $A$ as two independent risk factors for low birthweight.
    For simplicity, we assume that $P(W=1 \given \doit(S,A)) = 1_{S \vee A}$, that is, a baby has low birth weight iff the mother smokes or the baby has a genetic abnormality.
    Assume $P(M=1 \given \doit(S,A,W)) = m_0 e^{\mu_S S + \mu_A A + \mu_W W}$, which models $S,A,W$ as three independent risk factors of neonatal mortality.
    The story suggests $\omega_S > 0$, $\omega_A > 0$.
    Also $\mu_S > 0$, $\mu_A \gg 0$, $\mu_W > 0$.
    Finally, we will assume that $P(A=1)$ is very small.
    We can now look at the risk factor for smoking on mortality, for low birth-weight babies:
    \begin{equation*}\begin{split}
      \frac{P(M=1 \given \doit(S)=1,W=1)}{P(M=1 \given \doit(S)=0,W=1)} 
      &= \frac{\sum_a P(M=1 \given \doit(S)=1,W=1,A=a) P(A=a \given \doit(S=1),W=1)}{\sum_a P(M=1 \given \doit(S)=0,W=1,A=a) P(A=a \given \doit(S=0),W=1)} \\
      &= \frac{\sum_a m_0 e^{\mu_S} e^{\mu_W} e^{\mu_A a} P(A=a \given \doit(S=1),W=1)}{\sum_a m_0 e^{\mu_W} e^{\mu_A a} P(A=a \given \doit(S=0),W=1)} \\
      &= \frac{e^{\mu_S} (e^{\mu_A} P(A=1) + P(A=0))}{e^{\mu_A}} \\
      &= e^{\mu_S} (P(A=1) + e^{-\mu_A} P(A=0)) 
    \end{split}\end{equation*}
    We can now end up in the situation of the paradox if
      $$e^{\mu_S} (P(A=1) + e^{-\mu_A} P(A=0)) < 1 $$
    which can happen if the risk factor $e^{\mu_S}$ is not too big, while the incidence of genetic abnormalities is very small and the risk of dying from that is larger than that of the risk of dying from a smoking mother.
    }
\end{enumerate}
